\documentclass[11pt]{article}
\usepackage[T1]{fontenc}
\usepackage{lmodern}
\usepackage[margin=1in]{geometry}
\usepackage{amsmath,amssymb,amsthm,mathtools,mathrsfs}
\usepackage{microtype,booktabs,array,enumitem}
\usepackage[dvipsnames]{xcolor}
\usepackage[colorlinks=true,linkcolor=MidnightBlue,citecolor=MidnightBlue,urlcolor=MidnightBlue]{hyperref}
\usepackage{fancyhdr}
\hypersetup{pdftitle={Positive factorizations and supersymmetric extensions of finite-interval Weil forms},pdfauthor={Working draft prepared for Edward Baker}}
\pagestyle{fancy}
\fancyhf{}
\fancyhead[L]{\small Positive factorizations of Weil forms}
\fancyhead[R]{\small Working draft 0.2}
\fancyfoot[C]{\thepage}
\setlength{\headheight}{14pt}
\setlength{\emergencystretch}{2em}
\setlist{itemsep=3pt,topsep=5pt}
\numberwithin{equation}{section}
\newtheorem{theorem}{Theorem}[section]
\newtheorem{proposition}[theorem]{Proposition}
\newtheorem{lemma}[theorem]{Lemma}
\newtheorem{corollary}[theorem]{Corollary}
\theoremstyle{definition}
\newtheorem{definition}[theorem]{Definition}
\theoremstyle{remark}
\newtheorem{remark}[theorem]{Remark}
\newcommand{\R}{\mathbb R}
\newcommand{\C}{\mathbb C}
\newcommand{\HH}{\mathcal H}
\newcommand{\DD}{\mathcal D}
\newcommand{\norm}[1]{\lVert#1\rVert}
\newcommand{\ip}[2]{\langle#1,#2\rangle}
\newcommand{\Qgam}{Q^{\gamma}}
\newcommand{\EE}{\mathcal E}
\newcommand{\Acal}{\mathcal A}
\DeclareMathOperator{\Dom}{Dom}
\DeclareMathOperator{\supp}{supp}
\DeclareMathOperator{\Rea}{Re}
\DeclareMathOperator{\diag}{diag}
\DeclareMathOperator{\artanh}{artanh}
\title{Positive factorizations and supersymmetric extensions\\of finite-interval Weil forms}
\author{Working draft prepared for Edward Baker}
\date{Version 0.2 \quad 11 September 2026}

\begin{document}
\maketitle
\begin{abstract}
We investigate whether positive auxiliary systems can explain the sign of
finite-interval Weil forms arising from the shifted Riemann zeta function.
The objective is an independently specified square-root or bulk construction,
followed by an exact arithmetic identification. We give an explicit nonlocal
factorization of the prime-free form, uniform for total interval length
$L\leq1/4$ and shift $0\leq\omega\leq1/2$, with lower bound greater than
$0.099$. The gamma kinetic multiplier admits an exact realization by a
positive tower of local auxiliary fields and a closed graded supercharge.
We also prove restrictions on proposed completions. Scalar gauge changes
cannot remove the off-diagonal sign obstruction. More generally, positive
two-point square factors fail already at $L=2/3$, as established by a
comparison-form test with an exact rational certificate. Comparison
positivity survives elimination of internal nodes in finite pairwise
networks. A coherent multicomponent toy model escapes this restriction. Reflection
folding and an explicit cubic weight give a positive factor for the complete
central odd gamma sector for $L\leq7/10$, with floor $1/100$. An exact
linear-input test shows that the first prime can stabilize a negative gamma
direction and that an independent-edge remainder fails.
We identify additional boundary and first-prime matching constraints, and
show that no closable norm factor of the full Weil form exists on ordinary
whole-line $L^2$ with all smooth compactly supported inputs in its domain.
These results delimit a supersymmetric research program; they supply neither
an all-depth positivity theorem nor a proof of the Riemann hypothesis.
\end{abstract}

\noindent\textbf{Draft status.}
This manuscript consolidates three exploratory investigations. Proofs are
included, with reproducible exact rational checkers for the comparison
obstruction, the odd-sector potential bound, and the first-prime sign test.
There has been no independent specialist review or complete priority audit.
The graded constructions are linear operator models, not a constructed
interacting superspace field theory. The small-window result does not extend
the stronger finite-depth certificates already recorded in the project.

\smallskip
\noindent\textbf{Keywords.} Weil positivity; auxiliary fields; supersymmetric
quantum mechanics; nonlocal forms; Schur complements; comparison matrices.

\clearpage
\tableofcontents
\clearpage

\section{Motivation and the question of mechanism}
\label{sec:motivation}

\subsection{From higher-dimensional extensions to positive boundary systems}

The research program motivating this paper began with an attempt to extend
the completed xi function into higher-dimensional settings and constrain its
zeros by rotational symmetry. Investigations using slice-regular and related
extensions repeatedly encountered the same difficulty: the proposed
extension encoded information already present in the original complex
function. An additional coordinate did not, by itself, create an additional
constraint on the zeros. This is the history reported by the researcher,
not a general impossibility theorem about higher-dimensional approaches.

The search then moved toward physical models, including gauge theory and
ideas inspired by bulk--boundary correspondences. The enduring question was
whether an arithmetic boundary observable could be derived from an
independently positive system. Here a ``bulk'' may be a linear field system,
an auxiliary channel space, or a network. The term does not assert an AdS
geometry or a conformal field theory. A sufficiently explicit positive
linear model would already be useful if its boundary response were exactly
the required zeta object.

The shifted zeta program supplies concrete targets for this question:
causal transfers, finite-interval generators, contraction defects, and
string or spectral representations. Its existing finite-depth work proves
restricted inequalities by analytic estimates and certificates
\cite{BakerFirst,BakerWeil}. Those results provide benchmarks for an
explanatory construction. They do not establish positivity at every depth,
and a sequence of finite certificates is not itself a rule forcing the
next certificate to succeed.

The intended direction of explanation is
\[
\text{specified arithmetic data}
\ \longrightarrow\ \text{independently positive system}
\ \longrightarrow\ \text{exact boundary identity}.
\]
The global quantifiers needed for the Riemann hypothesis (RH) come only
after these steps and require a separate limiting or compatibility theorem.
Reconstructing a positive system from a measure already assumed positive
does not supply the missing direction of implication.

\subsection{Why a supersymmetric extension is a reasonable test}

Supersymmetric quantum mechanics suggests beginning with a defining operator
$A$ and a positive inner product. Its adjoint then yields a nonnegative
Hamiltonian $A^*A$, or a graded Hamiltonian with two partner blocks. The
substantive work is to define $A$ independently and identify the resulting
energy or observable with the arithmetic target.

Signed terms in a displayed operator do not rule out such a mechanism.
For example, on a compactly supported smooth core, a real smooth function
$W$ gives
\begin{equation}
(\partial_x+W)^*(\partial_x+W)
=-\partial_x^2+W^2-W'.
\end{equation}
The potential need not be pointwise positive, but the terms are constrained
by a common square. On an interval the same statement also requires the
correct domain and boundary terms. The present problem is to find an
analogous constraint relating the gamma, pole, and prime contributions.

The positive component-space interpretation of fermionic variables is
standard in supersymmetric quantum mechanics \cite{Singleton}. A formal
Berezin integral is not a positive measure, and a supertrace is a difference
of ordinary traces. Boson--fermion cancellation therefore does not establish
positivity of an arbitrary reduced observable. Likewise, nilpotence of a
BRST differential does not establish positivity on its physical quotient.
A BRST construction would require a specific gauge redundancy and a
separate argument about the induced inner product. No such gauge system is
assumed here.

\subsection{Results and limits of the present paper}

Theorem~\ref{thm:small} constructs a positive factor on a small prime-free
window. Theorem~\ref{thm:tower} realizes the gamma kinetic term through local
auxiliary fields; the reduction produces the correct nonlocal boundary
order. Theorem~\ref{thm:pair} excludes a larger class of pairwise factors,
and Proposition~\ref{prop:schur} explains why adding internal nodes to a
finite pairwise network does not automatically help.
Section~\ref{sec:coherent} exhibits a coherent positive model outside that
class and identifies two elementary completion failures.
Theorem~\ref{thm:oddcompletion} gives an explicit completion on the odd
subspace through the prime-free slab; Proposition~\ref{prop:primestabilizes}
certifies the stabilizing role of the first prime on a linear input.
Theorem~\ref{thm:global} addresses the input space for a global factor.

These statements have different scopes. A small-interval construction does
not prove all-depth positivity. A restriction on a class of factors does
not prove negativity of the Weil form. A coherent toy model establishes
possibility within a larger class, not an arithmetic selection principle.
The proposed continuation in Section~\ref{sec:program} is kept separate
from the proved statements.

There is substantial prior literature on Weil forms and Hilbert-space
realizations. In particular, Suzuki gives norm-identity criteria related to
RH \cite{SuzukiHilbert} and a screw-function treatment of finite-interval
operators and small-window arguments \cite{SuzukiScrew}. We do not claim that
the possibility of a norm formulation, a graded block construction, or a
small-window Dirichlet argument is new. The application-specific formulas
and restrictions below are recorded with proofs; their priority requires
further review.

\section{The forms, domains, and graded convention}
\label{sec:target}

Let $I_L=(-L/2,L/2)$, where $L$ is the \emph{total} interval length.
Extend $f\in L^2(I_L)$ by zero to $\widetilde f$ on $\R$, and use
\begin{equation}
\widehat f(\tau)=\int_{I_L}f(x)e^{-i\tau x}\,dx,
\qquad
\DD_{\log,L}=\left\{f:\int_\R\log(2+|\tau|)|\widehat f(\tau)|^2d\tau<\infty\right\}.
\end{equation}
We use inner products conjugate-linear in the first variable. Write
\begin{equation}
C(f)=\int_{I_L}f(x)\cosh(x/2)\,dx,
\qquad S(f)=\int_{I_L}f(x)\sinh(x/2)\,dx.
\end{equation}
For $a>0$, define the truncated translation
$(T_af)(x)=\widetilde f(x-a)$ for $x\in I_L$.

The central finite-interval form, in the normalization of
\cite{BakerWeil}, is
\begin{align}
Q_{0,L}[f]={}&\frac1{2\pi}\int_\R
\left(\Rea\psi(\tfrac14+\tfrac{i\tau}{2})-\log\pi\right)
|\widehat f(\tau)|^2\,d\tau\notag\\
&+2|C(f)|^2-2|S(f)|^2
-\sum_{\log n<L}\frac{\Lambda(n)}{\sqrt n}
\ip{f}{(T_{\log n}+T_{\log n}^*)f}.
\label{eq:weil}
\end{align}
Here $\psi=\Gamma'/\Gamma$, and $\Lambda$ is the von Mangoldt function.
The prime-free form $\Qgam_{0,L}$ omits the last sum and agrees with
$Q_{0,L}$ for $L\leq\log2$. The Fourier multiplier differs from
$\log(2+|\tau|)$ by a bounded function \cite{DLMF}; the other terms are
bounded on each fixed interval. Thus these forms are closed and semibounded
on $\DD_{\log,L}$, regardless of their sign. Smooth compactly supported
functions form a core: first shrink the support by dilation, then mollify;
the weighted Fourier norm converges in each step.

Set
\begin{align}
w_0&=\psi(1/4)-\log\pi
=-\gamma_E-\pi/2-3\log2-\log\pi,\label{eq:w0}\\
j(t)&=\frac{e^{-t/2}}{1-e^{-2t}},
& R(t)&=2\cosh(t/2)-j(t),\qquad t>0.\label{eq:jR}
\end{align}
The shifted gamma form can be defined without an endpoint ambiguity by
\begin{equation}
\Qgam_{\omega,L}[f]=\Qgam_{0,L}[f]
+\int_{I_L^2}(\cosh(\omega|x-y|)-1)R(|x-y|)
\overline{f(x)}f(y)\,dx\,dy,
\label{eq:shift}
\end{equation}
for $0\leq\omega\leq1/2$. This is the bounded perturbation identity proved
for the finite-horizon shift generator in \cite{BakerFirst}. The kernel is
continuous across the diagonal after assigning its limiting value there.
In particular, \eqref{eq:shift} includes the distributional correction
required at $\omega=1/2$ in a Fourier-plus-residue description.

\begin{lemma}[Graded realization]\label{lem:graded}
Let $A:\Dom A\subset\HH_0\to\HH_1$ be closed and densely defined between
positive Hilbert spaces. Then
\begin{equation}
\mathscr S=\begin{pmatrix}0&A^*\\A&0\end{pmatrix},
\quad \Dom\mathscr S=\Dom A\oplus\Dom A^*,
\end{equation}
is self-adjoint and
\begin{equation}
\mathscr S^2=\diag(A^*A,AA^*)\geq0.
\label{eq:graded}
\end{equation}
Equivalently, $q(f,g)=(0,Af)$, with domain $\Dom A\oplus\HH_1$,
is nilpotent and $\{q,q^*\}=\mathscr S^2$ on the natural operator domain.
\end{lemma}
\begin{proof}
Computing the adjoint of the off-diagonal block gives the same block and
domain because $A^{**}=A$. Squaring on its natural domain gives
\eqref{eq:graded}. The nilpotent formulation is the same calculation.
\end{proof}

This lemma organizes an independently constructed factor. Choosing
$A=Q^{1/2}$ for an operator whose sign is unknown would assume the premise
we seek to establish.

\section{An explicit positive factor on a small window}
\label{sec:small}

\subsection{The full-output jump identity}

The digamma difference integral \cite[\S5.9]{DLMF} gives
\begin{equation}
\Rea\psi(1/4+i\tau/2)-\log\pi
=w_0+2\int_0^\infty j(t)(1-\cos\tau t)\,dt.
\end{equation}
If $U_t$ denotes translation on $L^2(\R)$, Plancherel implies
\begin{equation}
\Qgam_{0,L}[f]=w_0\norm f^2
+\int_0^\infty j(t)\norm{\widetilde f-U_t\widetilde f}_{L^2(\R)}^2dt
+2|C(f)|^2-2|S(f)|^2.
\label{eq:jump}
\end{equation}
The norm includes output outside $I_L$. Dropping that output would change
the boundary contribution and invalidate the subsequent factorization.

Define
\begin{equation}
F(a)=\int_a^\infty j(t)dt
=\artanh(e^{-a/2})+\arctan(e^{-a/2}),\qquad a>0,
\end{equation}
and put $h(t)=-R(t)$. Splitting the jump integral into interior and exterior
pairs, then expanding the pole kernel, gives
\begin{align}
\Qgam_{0,L}[f]&=\frac12\int_{I_L^2}h(|x-y|)|f(x)-f(y)|^2dx\,dy
+\int_{I_L}\kappa_{0,L}(x)|f(x)|^2dx,\label{eq:diff}\\
\kappa_{0,L}(x)&=w_0+F(x+L/2)+F(L/2-x)
+8\sinh(L/4)\cosh(x/2).\label{eq:kappa}
\end{align}
The endpoint logarithms in $F$ are part of the form. They are integrable
against smooth functions. Identity~\eqref{eq:diff} remains valid when its
weights lose positivity.

\subsection{The positive range and shift-uniform estimate}

\begin{theorem}[Small-window factor]\label{thm:small}
For $0<L\leq1/4$ and $0\leq\omega\leq1/2$, define
\begin{align}
h_\omega(t)&=\cosh(\omega t)h(t),\\
\kappa_{\omega,L}(x)&=\kappa_{0,L}(x)
-\int_{I_L}(\cosh(\omega|x-y|)-1)h(|x-y|)dy.\label{eq:kappashift}
\end{align}
Both weights are nonnegative, and
\begin{equation}
\kappa_{\omega,L}(x)\geq c_*:=
\log(4/\pi)-\gamma_E+\frac7{16}
-\frac{e^{1/8}\cosh(1/8)}{512}>0.099.
\label{eq:cstar}
\end{equation}
The closed map from $L^2(I_L)$ to $L^2(I_L^2)\oplus L^2(I_L)$ given by
\begin{equation}
\Acal_{\omega,L}f=
\left(\sqrt{h_\omega(|x-y|)/2}\,[f(x)-f(y)],
\ \sqrt{\kappa_{\omega,L}(x)}f(x)\right)
\label{eq:smallA}
\end{equation}
has domain $\DD_{\log,L}$ and satisfies
\begin{equation}
\Qgam_{\omega,L}[f]=\norm{\Acal_{\omega,L}f}^2
\geq c_*\norm f^2.
\end{equation}
\end{theorem}

\begin{proof}
The sign of $R(t)$ is that of $q^3-q-1$ with $q=e^t$. Its unique positive
sign-change point is
\begin{equation}
t_0=\log\rho=0.2811995743\ldots,\qquad \rho^3=\rho+1.
\label{eq:tzero}
\end{equation}
In particular, $h(t)>0$ for $0<t\leq1/4$.

Since $j$ decreases, $F$ is convex. Both the sum of the two $F$ terms in
\eqref{eq:kappa} and its hyperbolic-cosine term attain their minimum at
$x=0$. Using elementary hyperbolic identities, this minimum is
\begin{equation}
-\gamma_E-\log(8\pi)+\log\coth(L/8)
-2\arctan(\tanh(L/8))+8\sinh(L/4).
\end{equation}
The inequalities $\coth u\geq1/u$, $\arctan(\tanh u)\leq u$, and
$\sinh u\geq u$ give
\begin{equation}
\kappa_{0,L}(x)\geq-\gamma_E-\log(\pi L)+7L/4
\geq\log(4/\pi)-\gamma_E+7/16=:c_0,
\label{eq:kappalower}
\end{equation}
where the last inequality uses monotonic decrease of the preceding
expression for $L\leq1/4$.

Expanding the difference squares shows that changing $h$ to $h_\omega$
produces the off-diagonal perturbation in \eqref{eq:shift}; the subtraction
in \eqref{eq:kappashift} cancels the extra diagonal. Moreover,
\[
0\leq h(t)\leq j(t)=\frac{e^{t/2}}{2\sinh t}\leq\frac{e^{1/8}}{2t},
\qquad
\cosh(\omega t)-1\leq\frac{t^2}{8}\cosh(1/8).
\]
Since $\int_{I_L}|x-y|dy\leq L^2/2$, the subtracted potential is at most
$e^{1/8}\cosh(1/8)/512$. This proves \eqref{eq:cstar}.

The norm identity follows first on the smooth core from \eqref{eq:diff}.
The corresponding graph norm is equivalent to the closed positive form
norm on $\DD_{\log,L}$. Closing the map on that core gives a closed
operator with this domain and preserves the identity. The sign estimate
precedes this closure argument; no unknown spectral square root is used.
\end{proof}

\subsection{Relation to the shift-energy identity}

On the smooth core, the finite-horizon transfer satisfies
$\partial_\omega V_{\omega,L}=-G_{\omega,L}V_{\omega,L}$, with symmetric
generator form $\Qgam_{\omega,L}$ in the prime-free region. Integration of
the energy derivative, followed by the form limit, gives
\begin{equation}
\norm f^2-\norm{V_{\omega,L}f}^2
=2\int_0^\omega
\norm{\Acal_{s,L}V_{s,L}f}^2ds,
\label{eq:radial}
\end{equation}
in the domain and limiting sense of \cite{BakerFirst}. Equation
\eqref{eq:radial} is a consequence of that transfer identity and the new
explicit factor, not a separately constructed bulk duality. Its range is
only a subset of the already certified prime-free slab.

\section{Two restrictions on direct scalar factors}
\label{sec:direct}

\begin{proposition}[Scalar conductances and scalar gauges]
For $L>t_0$, the gamma form cannot be written as nonnegative scalar
conductances multiplying $|f(x)-f(y)|^2$, plus an arbitrary multiplication
form. A positive scalar rescaling, or a single position-dependent phase
change of the input, does not remove this obstruction.
\end{proposition}
\begin{proof}
A nonnegative difference conductance has nonpositive off-diagonal kernel,
whereas $R(|x-y|)$ is positive on an open set when $L>t_0$.
A positive rescaling preserves these signs. For a phase change, choose
$t_0/2<d<\min(t_0,L/2)$ and three interior points with consecutive
separation $d$. The signs of the three kernel entries are $-,-,+$.
Their product is positive and is invariant under a scalar phase gauge.
Three negative real entries would have negative product. The argument can
be made on small separated neighborhoods, so it is an obstruction to an
almost-everywhere kernel identity, not merely to values at three points.
\end{proof}

Independent edge phases are more general than a scalar gauge. They are
addressed in Section~\ref{sec:comparison}; the preceding argument alone
does not exclude them.

\begin{proposition}[Local first-order boundary factors]
No finite-component map $Af=a(x)f'+b(x)f$, with locally bounded smooth
coefficients and the ordinary positive component $L^2$ norm, has
$\norm{Af}^2=Q_{0,L}[f]$ on all $C_c^\infty(I_L)$.
\end{proposition}
\begin{proof}
For fixed $\chi\in C_c^\infty(I_L)$, put $f_N=\chi e^{iNx}$.
The digamma asymptotic and the bounded remaining terms imply
\begin{equation}
Q_{0,L}[f_N]=\log N\,\norm\chi^2+O(1).
\end{equation}
For the multiplier term, translate the Fourier variable by $N$ and use
rapid decay of $\widehat\chi$ outside $|\tau-N|\leq N/2$.
If $a$ is nonzero, choose $\chi$ where $\norm{a\chi}>0$; then
$\norm{Af_N}^2=N^2\norm{a\chi}^2+O(N)$, a contradiction.
If $a$ vanishes identically, the left side remains bounded in $N$, also a
contradiction. The same argument applies to a finite vector of rows.
\end{proof}

This proposition concerns the square root acting directly on the original
boundary coordinate. Eliminating local auxiliary variables can change the
order of the reduced operator, as the next construction demonstrates.

\section{A positive auxiliary-field realization of the gamma kinetic term}
\label{sec:tower}

Define $a_k=2k+1/2$ and
\begin{equation}
B(s)=\sum_{k=0}^\infty\frac2{a_k}\frac{s}{a_k^2+s},\qquad s\geq0.
\label{eq:B}
\end{equation}
The geometric expansion $j(t)=\sum_{k\geq0}e^{-a_kt}$ and the digamma
difference integral give
\begin{equation}
B(\tau^2)=\Rea\psi(1/4+i\tau/2)-\psi(1/4).
\label{eq:Bdigamma}
\end{equation}
Consequently the nonnegative kinetic form is
\begin{equation}
K[f]=\frac1{2\pi}\int_\R B(\tau^2)|\widehat f(\tau)|^2d\tau.
\end{equation}

\begin{theorem}[Positive local tower]\label{thm:tower}
For $a>0$, let
\begin{equation}
\EE_a[f,u]=\frac2a\int_\R
\left(|f-u|^2+a^{-2}|u'|^2\right)dx,
\quad f\in L^2(\R),\quad u\in H^1(\R).
\label{eq:Ea}
\end{equation}
Then, with extended values allowed,
\begin{equation}
K[f]=\inf_{(u_k)}\sum_{k=0}^\infty\EE_{a_k}[f,u_k].
\label{eq:minK}
\end{equation}
For $K[f]<\infty$, the infimum is attained by
\begin{equation}
\widehat u_k(\tau)=\frac{a_k^2}{a_k^2+\tau^2}\widehat f(\tau).
\label{eq:uk}
\end{equation}
The unreduced energy is the norm square of a closed densely defined map
between positive component Hilbert spaces and therefore has the graded
realization of Lemma~\ref{lem:graded}.
\end{theorem}
\begin{proof}
At a single frequency, completion of the square gives
\begin{align}
|f-u|^2+\frac{\tau^2}{a^2}|u|^2
={}&\frac{\tau^2}{a^2+\tau^2}|f|^2\notag\\
&+\left(1+\frac{\tau^2}{a^2}\right)
\left|u-\frac{a^2}{a^2+\tau^2}f\right|^2.
\end{align}
Plancherel gives the minimum of \eqref{eq:Ea}. Summing and using monotone
convergence proves \eqref{eq:minK}. Each multiplier in \eqref{eq:uk}
maps $L^2$ into $H^1$, and the minimum has finite total energy whenever
$K[f]$ is finite.

For a precise infinite-component space, set
\begin{equation}
\HH_B=L^2(\R)\oplus\ell^2\bigl(L^2(\R);2/a_k^3\bigr),
\qquad
\HH_F=\ell^2\bigl(L^2(\R)\oplus L^2(\R)\bigr),
\end{equation}
and define
\begin{equation}
A(f,(u_k))=
\left(\sqrt{2/a_k}(f-u_k),\sqrt{2/a_k^3}\,u_k'\right)_{k\geq0}.
\label{eq:towerA}
\end{equation}
Give $A$ its maximal domain: all $u_k\in H^1$ for which the displayed
sequence belongs to $\HH_F$. It is closed. Indeed, convergence in
$\HH_B$ gives convergence of every component in $L^2$, and convergence
of $A$ gives convergence of every weak derivative; closedness of weak
differentiation identifies the limit and its residuals.

The domain is dense. Choose smooth $f$, let finitely many smooth $u_k$
be independent, and set $u_k=f$ in the tail. The tail derivative energy
is finite since $\sum a_k^{-3}<\infty$. Such vectors are dense in the
weighted product space by truncation and approximation. Finally, the
minimizers satisfy $\norm{u_k}\leq\norm f$, so their auxiliary Hilbert
norm is finite; their derivative and difference energies sum to $K[f]$.
Lemma~\ref{lem:graded} now applies to \eqref{eq:towerA}.
\end{proof}

The fields $u_k$ live on the whole line even when $f$ is supported in
$I_L$. Truncating them at the input endpoints changes the minimization.
The square-root map uses only first derivatives and multiplication in the
unreduced coordinate. Its reduced boundary multiplier is logarithmic,
which explains precisely why the direct order argument of
Section~\ref{sec:direct} does not apply.

\begin{corollary}[An extra-coordinate extension]
The function $B$ in \eqref{eq:B} is a complete Bernstein function. Hence
$B(-\partial_x^2)$ has a positive extension realization of the type
established by Kwa\'snicki and Mucha, allowing a measure-valued coefficient.
\end{corollary}
\begin{proof}
The representation
\[
B(s)=\int_0^\infty\frac{s}{s+r}\,\nu(dr),
\qquad \nu=\sum_{k\geq0}\frac2{a_k}\delta_{a_k^2},
\]
has a positive measure satisfying
$\int(1+r)^{-1}\nu(dr)<\infty$. This is the complete Bernstein
representation, and the extension conclusion follows from
\cite{KwasnickiMucha}, including its general measure-coefficient theorem.
\end{proof}

The extension coefficient is not computed here. This corollary concerns
the explicitly positive function $B$, not the shifted zeta Weyl function.
Moreover, the exact gamma form still reads
\begin{equation}
\Qgam_{0,L}[f]=K[\widetilde f]+w_0\norm f^2
+2|C(f)|^2-2|S(f)|^2.
\label{eq:remaining}
\end{equation}
The negative normalization and odd pole contribution remain. They cannot
be removed by changing the zero of energy without changing the target.
The positive tower is an ingredient for a completion, not the completion.
Related positive diffusive gamma expansions already appear in
\cite{BakerFirst}; the point here is the explicit variational and graded
realization at the generator level.

\section{A comparison-form obstruction to pairwise factors}
\label{sec:comparison}

\subsection{The class of factors and its necessary condition}

The scalar gauge obstruction does not control arbitrary phases attached
to edges. A stronger test uses the diagonal contributions forced by each
positive two-point square.

\begin{definition}[Positive pairwise representation]\label{def:pair}
A positive pairwise representation on $I$ is a sum or integral of terms
\begin{equation}
\norm{u(x,y)f(x)+v(x,y)f(y)}^2,
\qquad x<y,
\label{eq:pairrow}
\end{equation}
plus a nonnegative multiplication form. The coefficient vectors may have
arbitrary positive Hilbert component spaces. After grouping channels for
each pair, their Gram coefficients are locally integrable away from
$x=y$. The nonnegative integral is finite on the smooth test core and
equals the proposed form there. Negative renormalizing counterterms are
not permitted as additional energy terms in this definition.
\end{definition}

For a fixed pair the Gram matrix has diagonal entries
$\alpha=\norm u^2$, $\beta=\norm v^2$, and off-diagonal entry
$c=\ip uv$. Cauchy--Schwarz gives $|c|^2\leq\alpha\beta$.
Replacing $c$ by $-|c|$ therefore preserves positivity of this
two-by-two matrix. Different endpoint magnitudes and independent edge
phases are all allowed in Definition~\ref{def:pair}.

\begin{lemma}[Necessary comparison positivity]\label{lem:comparison}
If $\Qgam_{0,L}$ has a positive pairwise representation, then
\begin{equation}
Q_{\mathrm{cmp},L}[f]:=\Qgam_{0,L}[f]
-2\int_{I_L^2}R_+(|x-y|)\overline{f(x)}f(y)dx\,dy
\label{eq:comparison}
\end{equation}
is nonnegative on nonnegative real smooth supported inputs. Here
$R_+=\max(R,0)$.
\end{lemma}
\begin{proof}
Polarization with disjointly supported tests identifies the grouped
off-diagonal coefficient with $R(|x-y|)$ almost everywhere away from
the diagonal. For nonnegative real inputs, replacing that coefficient by
$-|R|$ subtracts exactly the second term in \eqref{eq:comparison}.
Each grouped pair remains positive by the preceding Gram-matrix argument,
and the nonnegative multiplication term is unchanged. The subtraction
has a bounded kernel, so the argument does not split the singular
diagonal into separately divergent integrals.
\end{proof}

\subsection{A counterexample at total length \texorpdfstring{$2/3$}{2/3}}

\begin{theorem}[Failure of positive pairwise factors]\label{thm:pair}
The form $\Qgam_{0,2/3}$ has no representation of
Definition~\ref{def:pair}. In particular, independent edge phases,
unequal endpoint weights, and arbitrary separate vector channels for
each pair cannot provide such a factorization.
\end{theorem}

\begin{proof}
Take $f=L^{-1/2}\mathbf1_{I_L}$. It belongs to the logarithmic form
domain. The jump identity and cancellation of the exponential tail give
\begin{equation}
\Qgam_{0,L}[f]=w_0+\frac2L\sum_{k=0}^\infty
\frac{1-e^{-a_kL}}{a_k^2}
+\frac{32}{L}\sinh^2(L/4),
\label{eq:constantgamma}
\end{equation}
while
\begin{equation}
Q_{\mathrm{cmp},L}[f]
=\Qgam_{0,L}[f]-\frac4L\int_0^L(L-t)R_+(t)dt.
\label{eq:constantcmp}
\end{equation}
For $L=2/3$ and $b=1/3$, the right side is at most
\begin{equation}
\mathcal U:=\Qgam_{0,2/3}[f]
-6\int_{1/3}^{2/3}(2/3-t)R(t)dt.
\label{eq:U}
\end{equation}
The exact rational enclosure established in
Appendix~\ref{app:certificate} proves $\mathcal U<-1/25$.

Choose nonnegative smooth interior cutoffs converging to the constant.
Each endpoint error is a scaled fixed profile of width $\varepsilon$;
changing variables in its weighted Fourier norm gives an upper bound
$O(\varepsilon(1+\log(1/\varepsilon)))$. Thus the cutoffs converge in
the logarithmic form norm. The bounded-kernel term in
\eqref{eq:comparison} also converges. A sufficiently close smooth cutoff
therefore has negative comparison form, contradicting
Lemma~\ref{lem:comparison}.
\end{proof}

For orientation, separate quadratures give
\[
\Qgam_{0,2/3}[f]\simeq0.17402032645,
\qquad Q_{\mathrm{cmp},2/3}[f]\simeq-0.06189510026.
\]
These decimal values are diagnostics; the proof uses the weaker exact
bound on $\mathcal U$. Since $2/3<\log2$, the obstruction occurs before
any prime is active. It concerns a class of positive representations,
not negativity of the gamma form itself.

\subsection{Eliminating hidden nodes in finite networks}

For a Hermitian matrix $G$, let $M(G)$ have the same diagonal and
off-diagonal entries $-|G_{ij}|$.

\begin{proposition}[Comparison positivity survives a Schur reduction]
\label{prop:schur}
If $M(G)\succeq0$, then eliminating internal nodes by positive-pivot
Schur complements preserves comparison positivity. In particular, every
finite positive network assembled from two-point squares and nonnegative
one-point terms has a boundary Schur complement $S$ with $M(S)\succeq0$.
\end{proposition}
\begin{proof}
Eliminate one node with diagonal $d>0$. The resulting entries are
$S_{ij}=G_{ij}-G_{i0}G_{0j}/d$. The diagonal of $M(S)$ equals that of
the Schur complement $N$ of $M(G)$, and, for $i\ne j$,
\begin{equation}
-|S_{ij}|\geq-|G_{ij}|-|G_{i0}||G_{0j}|/d=N_{ij}.
\end{equation}
Both $M(S)$ and $N$ have nonpositive off-diagonal entries. For every
complex vector $z$, writing $|z|$ componentwise gives
\begin{equation}
z^*M(S)z\geq |z|^{\mathsf T}M(S)|z|
\geq |z|^{\mathsf T}N|z|\geq0.
\end{equation}
The last step follows from positivity of the Schur complement of
$M(G)$. A zero diagonal in a positive semidefinite comparison matrix
has a zero row and column, so such nodes decouple. Iterate the argument.
Finally, the grouped two-by-two Gram argument proves $M(G)\succeq0$
for a positive pairwise network.
\end{proof}

The proposition is finite dimensional. It passes to a limit if the
comparison forms themselves converge on the relevant tests. No claim is
made here that every singular continuum extension admits that passage.
The theorem does establish that adding arbitrarily many nodes to a
finite scalar magnetic spring network does not, by itself, evade the
comparison restriction.

\section{Coherent couplings and two completion constraints}
\label{sec:coherent}

\subsection{A positive graded model outside the pairwise class}

Consider three complex boundary amplitudes and one auxiliary amplitude:
\begin{equation}
\EE(f,u)=|f_1+f_2-u|^2+|f_3+u|^2.
\label{eq:coherent}
\end{equation}
Completion of the square gives
\begin{equation}
\inf_u\EE(f,u)=\tfrac12|f_1+f_2+f_3|^2,
\qquad u=\tfrac12(f_1+f_2-f_3).
\end{equation}
The boundary matrix is half the all-ones matrix; its eigenvalues are
$0,0,3/2$. Its comparison matrix has eigenvalues $-1/2,1,1$.
Thus a positive bulk can produce a boundary form outside the pairwise
comparison class.

The independently specified factor is
\begin{equation}
A=\begin{pmatrix}1&1&0&-1\\0&0&1&1\end{pmatrix}.
\end{equation}
Lemma~\ref{lem:graded} supplies its graded realization. The first row
mixes three scalar amplitudes before taking the norm; it is not a
two-endpoint spring. Dividing it into independent pairwise energies would
discard the cross terms responsible for the example.

This is an architectural example, not a zeta model. Its positivity is
already present before the grading is introduced. A useful supersymmetric
completion of the arithmetic form would need a principle that forces
the analogous coherent couplings and their normalization.

\subsection{The pole term and a growing Green function}

The identity
\begin{equation}
R(t)=e^{t/2}-\sum_{k\geq1}e^{-(2k+1/2)t}
\label{eq:Rsplit}
\end{equation}
suggests replacing one decaying massive propagator by a growing one.
The following elementary computation identifies the boundary cost of
this particular attempt.

\begin{proposition}[Growing propagator boundary defect]
For $a>0$, the kernel
\begin{equation}
G_a(x,y)=-\frac1{2a}e^{a|x-y|},\qquad 0<x,y<L,
\end{equation}
is the Green function of $-\partial_x^2+a^2$ with boundary conditions
$u'(0)=-au(0)$ and $u'(L)=au(L)$. Its quadratic form is not nonnegative.
\end{proposition}
\begin{proof}
The derivative jump at $x=y$ and the two endpoint derivatives verify
the Green equation and Robin conditions. The homogeneous equation has
no nonzero solution satisfying both conditions, so the inverse is
well defined. Integration by parts gives the form on $H^1(0,L)$,
\begin{align}
\EE_R[u]&=\int_0^L(|u'|^2+a^2|u|^2)dx
-a|u(0)|^2-a|u(L)|^2\notag\\
&=\norm{u'-au}^2-2a|u(0)|^2.
\end{align}
The form-domain test $u(x)=e^{ax}$ gives $\EE_R[u]=-2a<0$.
\end{proof}

This calculation excludes treating the replacement channel as independently
positive. It does not exclude stabilization through a larger coherent
system. It also illustrates why a formal differential square without its
boundary domain is insufficient.

\subsection{The first prime forces joint diagonal accounting}

Translate the interval to $(0,L)$ and assume $\log2<L\leq\log3$.
Let $a=\log2$, $c=\log2/\sqrt2$, and $T=T_a$. Since $L<2a$,
$T^2=0$. Its initial and final projections are
\begin{equation}
P=T^*T=\mathbf1_{(0,L-a)},\qquad
F=TT^*=\mathbf1_{(a,L)}.
\end{equation}
The first-prime correction is $-c(T+T^*)$. On each paired component it
has eigenvalues $\pm c$, so it cannot be a separate positive additive
channel.

The explicit edge factor $B_p=\sqrt c(P-T^*)$ satisfies
\begin{equation}
B_p^*B_p=c(P+F-T-T^*).
\end{equation}
Hence the exact identity is
\begin{equation}
Q_{0,L}[f]=\norm{B_pf}^2+
\left(\Qgam_{0,L}[f]-c\ip f{(P+F)f}\right).
\label{eq:primebudget}
\end{equation}
Proposition~\ref{prop:primestabilizes} shows that the remainder fails
to be nonnegative throughout this interval. At positive shift the
arithmetic coefficient becomes $c\cosh(\omega\log2)$; the same accounting
applies to that correction.

There is a complementary restriction on an auxiliary-field approach.
For a quadratic block with a fixed bare boundary block $A_0$ and positive
invertible auxiliary block $C_0$, minimization gives
\begin{equation}
\begin{pmatrix}A_0&B_0^*\\B_0&C_0\end{pmatrix}
\quad\longmapsto\quad A_0-B_0^*C_0^{-1}B_0.
\label{eq:schurminus}
\end{equation}
The change is negative semidefinite, whereas the first-prime correction
is indefinite. Thus adding an auxiliary Gaussian field while leaving
the bare boundary block fixed cannot reproduce that correction either.
Equations~\eqref{eq:primebudget} and \eqref{eq:schurminus} require a
successful model to change the joint couplings and diagonal terms together.

\section{Reflection folding and an explicit odd-sector completion}
\label{sec:oddcompletion}

The comparison obstruction leaves open the possibility of mixing amplitudes
before applying a positive factor. Reflection gives one such mixing rule
already specified by the target. This section completes the central
prime-free form on the odd subspace. It does not complete the even subspace.

Put $\ell=L/2$. For $\varepsilon\in\{+1,-1\}$, the map
$f\mapsto g=\sqrt2 f|_{(0,\ell)}$ is unitary from the parity subspace
$f(-x)=\varepsilon f(x)$ onto $L^2(0,\ell)$. The folded off-diagonal
kernel of $\Qgam_{0,L}$ is
\begin{equation}
R(|x-y|)+\varepsilon R(x+y),\qquad 0<x,y<\ell.
\label{eq:foldedkernel}
\end{equation}
Since $j'<0$, we have $R'(t)=\sinh(t/2)-j'(t)>0$. Thus the odd
folded kernel is negative. Write its positive conductance as
\begin{equation}
h_-(x,y)=R(x+y)-R(|x-y|)>0\quad(x\ne y).
\end{equation}
In contrast, for $L>2t_0$ the even folded kernel is positive on an open
set: take $t_0<x<\ell$ and $y$ sufficiently small. The same positive
conductance construction therefore does not cover the even sector at
$L=\log2$.

\begin{theorem}[Explicit odd gamma factor]\label{thm:oddcompletion}
For $0<L\leq7/10$, set
\begin{equation}
\phi_L(x)=x\left(1-\frac{4x^2}{5\ell^2}\right),\quad 0<x<\ell.
\label{eq:cubicphi}
\end{equation}
Let $v_L=(Q^-_{\gamma,L}\phi_L)/\phi_L$, where $Q^-_{\gamma,L}$
is the folded odd operator. Then $v_L(x)>1/100$ throughout $(0,\ell)$,
and for every odd $f\in\DD_{\log,L}$, with $g=\sqrt2 f|_{(0,\ell)}$,
\begin{align}
\Qgam_{0,L}[f]={}&\frac12\int_{(0,\ell)^2}
 h_-(x,y)\phi_L(x)\phi_L(y)
 \left|\frac{g(x)}{\phi_L(x)}-\frac{g(y)}{\phi_L(y)}\right|^2dx\,dy
 \notag\\
&+\int_0^\ell v_L(x)|g(x)|^2dx
\ \geq\frac1{100}\norm f^2.
\label{eq:oddfactor}
\end{align}
The two displayed terms give an explicit closed norm factor and hence a
graded realization. For $L\leq\log2$ this is the full central Weil form
on odd inputs. Between $\log2$ and $7/10$ the statement concerns the gamma
form only.
\end{theorem}

\begin{proof}
We prove the sign of $v_L$ directly from the cubic, without a ground-state
or spectral positivity assumption. Set
\[
r(t)=j(t)-\frac1{2t},\qquad b=\frac4{5\ell^2},\qquad
S_\phi=\int_{-\ell}^{\ell}(y-by^3)\sinh(y/2)dy.
\]
The full-output jump identity, or direct integration of its polynomial
input, gives the following pointwise formula for $0<x<\ell$:
\begin{align}
v_L(x)={}&-\gamma_E-\log(2\pi)-\tfrac12\log(\ell^2-x^2)
 +\frac{1+b\ell^2/2-11bx^2/6}{1-bx^2}\notag\\
&-\frac{\displaystyle\int_0^\ell
 [r(|x-y|)-r(x+y)]\phi_L(y)dy
       +2S_\phi\sinh(x/2)}{\phi_L(x)}.
\label{eq:oddpotential}
\end{align}
For clarity about the normalization, one uses
\[
F(t)=\log2+\pi/4-\tfrac12\log t-\int_0^t r(s)ds.
\]
The constant part of $r$ integrates to zero against the odd polynomial.
The singular $1/(2t)$ jump contributes
$x-b(11x^3/6-\ell^2x/2)$; these facts give
\eqref{eq:oddpotential}, including its pole term.

We need only an elementary bound for $r'$. With $u=t/2$,
\begin{equation}
j(t)=\tfrac14\bigl(\operatorname{csch}u+\operatorname{sech}u\bigr),
\qquad
-r'(t)=\tfrac18\left(\frac{\cosh u}{\sinh^2u}-\frac1{u^2}
                   +\operatorname{sech}u\tanh u\right).
\end{equation}
For $0<u\leq7/20$,
\begin{equation}
0\leq\frac{\cosh u}{\sinh^2u}-\frac1{u^2}\leq\frac16.
\label{eq:cschbound}
\end{equation}
For the upper bound, expand
$(1+u^2/6)\sinh^2u-u^2\cosh u$ in powers: its coefficients vanish
through degree four and are nonnegative thereafter. For the lower bound,
use $\sinh u/u\leq(1-u^2/6)^{-1}$ and $\cosh u\geq1+u^2/2$.
The inequality $(1+v/2)(1-v/6)^2\geq1$ for $0\leq v\leq1$ finishes
the comparison. Therefore
\begin{equation}
0\leq-r'(t)\leq D:=31/480\qquad(0<t\leq7/10).
\end{equation}
Put $z=x^2/\ell^2$. The mean-value bound and an elementary integration give
\begin{align}
0\leq\frac{\int_0^\ell[r(|x-y|)-r(x+y)]\phi_L(y)dy}{\phi_L(x)}
&\leq\frac{2D\ell^2(3/10-z/6+z^2/25)}{1-4z/5},\\
\frac{2S_\phi\sinh(x/2)}{\phi_L(x)}
&\leq\frac{(13/75)\ell^3\cosh^2(\ell/2)}{1-4z/5}.
\end{align}
Here $\int_0^\ell\min(x,y)\phi_L(y)dy
=x\ell^2(3/10-z/6+z^2/25)$, and
$S_\phi\leq(13/75)\ell^3\cosh(\ell/2)$.

Let $\ell_*=7/20$. Since all loss bounds increase with $\ell$ and
$\cosh^2(7/40)<33/32$, \eqref{eq:oddpotential} is bounded below by
\begin{align}
\mathcal V(z)={}&-\gamma_E-\log(7\pi/10)-\tfrac12\log(1-z)\notag\\
&+\frac{\displaystyle 7/5-22z/15
-2D\ell_*^2(3/10-z/6+z^2/25)
-(13/75)\ell_*^3(33/32)}{1-4z/5}.
\label{eq:oddscalar}
\end{align}
The exact rational interval calculation described in
Appendix~\ref{app:round3} proves $\mathcal V(z)>1/100$ for all
$0\leq z<1$. Thus the sign proof is a scalar estimate for an explicitly
chosen polynomial, not a spectral certificate for the target operator.

Expanding the difference square in \eqref{eq:oddfactor} now proves the
identity on the odd smooth core, using \eqref{eq:foldedkernel} and the
definition of $v_L$. One can first cut away the diagonal and then take the
form limit. The cubic is positive on $(0,\ell)$; at zero, smooth odd
inputs have a bounded quotient by it. The graph norm of the displayed
factor equals the norm of the closed positive form. Closing from the
core gives the stated domain and identity, and Lemma~\ref{lem:graded}
applies.
\end{proof}

Reflection is an independently specified coherent operation here: in the
original input variables the odd projection uses $f(x)-f(-x)$ before
forming differences between two positive coordinates. This involves up to
four original amplitudes. It is not a representation of the full form by
the two-point factors excluded in Theorem~\ref{thm:pair}. The polynomial
choice is an explicit trial weight, not an arithmetic selection principle
for the even sector or for arbitrary depth. No novelty is claimed for the
general ground-state transform used in the proof.

\subsection{The first prime can supply essential stabilization}

The diagonal accounting in \eqref{eq:primebudget} has a concrete failure
within the first-prime interval.

\begin{proposition}[A certified linear-input test]\label{prop:primestabilizes}
At $L=1$, put $f(x)=\sqrt{12}\,x$ on $(-1/2,1/2)$, so $\norm f=1$.
With $P,F,c$ as in \eqref{eq:primebudget},
\begin{equation}
\Qgam_{0,1}[f]<-0.12,\qquad Q_{0,1}[f]>0.26,\qquad
\Qgam_{0,1}[f]-c\ip f{(P+F)f}<-0.58.
\label{eq:primesigns}
\end{equation}
Consequently the proposed remainder in \eqref{eq:primebudget} is not
nonnegative throughout the first-prime interval.
\end{proposition}
\begin{proof}
For the unit linear input on a general interval, the autocorrelation is
$r_f(t)=1-3t/L+2t^3/L^3$ for $0\leq t\leq L$. At $L=1$ only $n=2$
is active. Writing $a=\log2$ gives the exact expressions
\[
Q_{0,1}[f]=\Qgam_{0,1}[f]-2c(1-3a+2a^3),\qquad
c\ip f{(P+F)f}=c\bigl[1-(2a-1)^3\bigr].
\]
The jump formula evaluates $\Qgam_{0,1}[f]$ by a convergent exponential
mass sum. Appendix~\ref{app:round3} gives that sum and proved tails.
Exact rational interval evaluation yields
\begin{align*}
\Qgam_{0,1}[f]&\in[-0.137027,-0.128850],\\
Q_{0,1}[f]&\in[0.268205,0.276382],\\
\Qgam_{0,1}[f]-c\ip f{(P+F)f}&\in[-0.598903,-0.590726].
\end{align*}
These displayed intervals are rounded outward; the checker asserts the
stronger rational inequalities \eqref{eq:primesigns}. The linear input is
in the logarithmic form domain. Core approximation preserves the strict
signs, so endpoint discontinuities do not create the obstruction.
\end{proof}

This proposition proves signs on a specified input, not positivity of the
full first-prime operator on every input. In particular, the gamma piece
itself cannot be an independently positive additive subsystem throughout
that interval: the arithmetic coupling rescues this odd direction.

\subsection{Which reflected channel the prime stabilizes}

For $\log2<L\leq\log3$, still write $a=\log2$ and $\ell=L/2$. On
$L^2(0,\ell)$ let $J=(a-\ell,\ell)$, let $P_J$ multiply by
$\mathbf1_J$, and define the partial reflection
\begin{equation}
(\mathcal R_a g)(x)=\mathbf1_J(x)g(a-x),\qquad
\Pi_\pm=\tfrac12(P_J\pm\mathcal R_a).
\end{equation}
The operators $\Pi_\pm$ are orthogonal projections onto the symmetric
and antisymmetric components of the cap $J$. Directly folding the
translation gives, for an input of parity $\varepsilon$,
\begin{equation}
-c\ip f{(T_a+T_a^*)f}
=-\varepsilon c\ip g{\mathcal R_a g}
=-\varepsilon c\bigl(\norm{\Pi_+g}^2-\norm{\Pi_-g}^2\bigr).
\label{eq:primefold}
\end{equation}
Thus the prime adds $c$ to the symmetric cap channel of the odd sector
and subtracts $c$ from its antisymmetric channel. The signs reverse in
the even sector. In these coordinates the independent-edge square is
$2c\norm{\Pi_{-\varepsilon}g}^2$, and the remainder subtracts
$c\norm{P_Jg}^2$. This exhibits precisely where that allocation loses
positivity. Equation~\eqref{eq:primefold} is an exact coherent matching
rule, but one of its two channels is still negative. A successful
completion must control both against the gamma part jointly.

\section{The input space of a global norm identity}
\label{sec:global}

Let $Q_W$ denote the full central Weil form on $C_c^\infty(\R)$ with no
fixed support bound. The following restriction does not concern the
finite-interval closed forms.

\begin{theorem}[No closable factor on ordinary whole-line $L^2$]
\label{thm:global}
There is no closable linear map $A:L^2(\R)\supset\Dom A\to\mathcal K$
to a positive Hilbert space, with $C_c^\infty(\R)\subset\Dom A$, such that
\begin{equation}
\norm{Af}_{\mathcal K}^2=Q_W[f]
\quad\text{for every }f\in C_c^\infty(\R).
\end{equation}
\end{theorem}
\begin{proof}
Such an identity implies Weil positivity and hence RH; we use the standard
Weil criterion as presented in \cite{SuzukiHilbert,SuzukiScrew}.
Under RH the explicit formula becomes
\begin{equation}
Q_W[f]=\sum_\gamma m_\gamma|\widehat f(\gamma)|^2,
\label{eq:sampling}
\end{equation}
where distinct zero ordinates are indexed by $\gamma$ and $m_\gamma$
is their multiplicity. Fix an ordinate $\gamma_0$ and choose
$\phi\in C_c^\infty(\R)$ with $\widehat\phi(0)=1$. Set
\begin{equation}
f_N(x)=N^{-1}e^{i\gamma_0x}\phi(x/N).
\end{equation}
Then $\norm{f_N}^2=N^{-1}\norm\phi^2\to0$, while
$\widehat f_N(\gamma)=\widehat\phi(N(\gamma-\gamma_0))$.
These sampling vectors converge in $\ell^2(m_\gamma)$ to the vector
equal to one at $\gamma_0$ and zero elsewhere. To justify norm
convergence, use isolation of $\gamma_0$, rapid Fourier decay, and the
zero-counting estimate $N_\zeta(T)=O(T\log T)$ to obtain a summable
dominating sequence, uniform for $N\geq1$.

Consequently $Q_W[f_N-f_M]\to0$, but $Q_W[f_N]\to m_{\gamma_0}>0$.
The assumed norm identity makes $Af_N$ a Cauchy sequence with nonzero
limiting norm while $f_N\to0$. This contradicts closability.
\end{proof}

The argument is an unconditional impossibility proof: assuming the factor
first supplies RH and then yields a contradiction. It neither assumes RH
as a result of this investigation nor challenges the possibility of norm
identities on other input spaces. Related measure-theoretic closability
questions occur in Wiener--Hopf form theory \cite{Yafaev}; the proof here
is direct for the present sampling measure.

\begin{corollary}
A bounded embedding $J:L^2(\R)\to\HH_B$ followed by a closed supercharge
component $S$ cannot give such a global factor if
$C_c^\infty(\R)\subset\Dom(SJ)$.
\end{corollary}
\begin{proof}
The composition $SJ$ is closed: if $f_n\to f$ and $SJf_n\to g$, then
$Jf_n\to Jf$, and closedness of $S$ applies. Theorem~\ref{thm:global}
therefore excludes the claimed norm identity.
\end{proof}

There is no contradiction with the kinetic tower, which realizes a
different form with a continuous Fourier density. There is also no
contradiction with finite-interval factors: the supports of $f_N$ grow
without bound. Test-function topologies, other independently defined
positive spaces, and unbounded boundary observables remain possible.
Adding a fixed finite number of ordinary Sobolev derivatives to the
whole-line norm does not remove this particular sequence: its derivative
norms still tend to zero.

\section{What a useful supersymmetric completion must establish}
\label{sec:program}

\subsection{The boundary object must be specified}

Three possible constructions need different proofs. A direct factor
requires $Q[f]=\norm{Af}^2$ on its specified input space. A minimized
bulk energy requires an identity with an infimum over auxiliary fields,
including the correct endpoint conditions. A spectral response requires
identification of a resolvent or covariance with the zeta function of
interest. Positivity of a Hamiltonian does not identify these objects
with one another.

For example, $H\geq0$ and a Hilbert-space vector $b$ imply
\[
\ip b{(H-z)^{-1}b}=\int_0^\infty\frac{d\mu_b(\lambda)}{\lambda-z},
\qquad d\mu_b\geq0.
\]
This is a useful model of a positive response, but an actual boundary
trace can be unbounded and can require a different formulation. To use
such a mechanism for zeta, the exact response and normalization must be
derived from independently specified dynamics.

The shifted transfer family has the form
\begin{equation}
\Theta_\omega(z)=
\frac{\xi(1/2-\omega-iz)}{\xi(1/2+\omega-iz)},
\qquad
\xi(s)=\tfrac12s(s-1)\pi^{-s/2}\Gamma(s/2)\zeta(s).
\end{equation}
Increasing support length adds logarithmic arithmetic delays; changing
$\omega$ changes the shifted family. Neither operation is automatically
the extra coordinate of an extension problem. The RG and bulk--boundary
language remains motivation until a theorem identifies the corresponding
dynamics and observables.

\subsection{Two concrete milestones}

The first milestone is a coherent positive model reproducing the complete
prime-free form throughout the known positive slab, with the fixed
normalization and pole terms in \eqref{eq:remaining}. The kinetic tower is
a controlled ingredient, and the comparison theorem excludes a broad
family of simple completions. The odd central sector now has the explicit
factor of Theorem~\ref{thm:oddcompletion}; the even sector and a comparable
shift-uniform completion remain to be constructed. A proposed model should be tested against
both its off-diagonal kernel and its diagonal terms.

The second milestone is the first prime. The joint model must derive the
indefinite correction and the diagonal accounting in
\eqref{eq:primebudget}; appending an independent positive square or using
\eqref{eq:schurminus} with the original boundary block fixed is
insufficient. Proposition~\ref{prop:primestabilizes} shows that the
independent-edge remainder actually fails, while \eqref{eq:primefold}
identifies the channels requiring joint control. Beyond the first prime,
overlapping delays impose further
compatibility conditions not investigated here.

An arithmetic or geometric rule forcing coherent mixing would be
substantive progress. A Ward identity could contribute if the model,
domain, and invariant positive inner product were independently defined
and the identity forced the required equality. At present, no such rule
has been found. A fitted positive matrix, a spectral square root of the
target, or an inverse realization conditioned on positivity does not
provide this missing explanation.

\subsection{What should be retained if this route stops}

The earlier higher-dimensional investigations warn that additional
coordinates can reproduce existing information without strengthening it.
The present calculations refine that lesson. Auxiliary dynamics can
change the mathematical class of a boundary operator, as demonstrated
by the gamma tower. Coherent mixing can change the class of positive
boundary couplings, as demonstrated by \eqref{eq:coherent}. Both are
real possibilities, but neither selects the arithmetic data needed for RH.

The obstructions should be used at their stated scope. A failure of
pairwise squares motivates testing coherent amplitudes; it does not
establish a successful coherent completion. A domain obstruction
requires a different boundary topology; merely renaming the original
$L^2$ space as superspace does not change it. These are design constraints
for future explorations rather than evidence that a global proof is near.

\appendix
\section{Exact rational certificate for the comparison test}
\label{app:certificate}

The sign used in Theorem~\ref{thm:pair} is verified with rational interval
arithmetic in the companion file \texttt{check\_round2.py}. This appendix
states the finite calculation and all truncation bounds.

Let $L=2/3$, $b=1/3$, $N=32$, and $a_k=2k+1/2$. Define
\begin{equation}
J(\lambda)=\int_b^L(L-t)e^{-\lambda t}dt
=\frac{e^{-\lambda b}[\lambda(L-b)-1]+e^{-\lambda L}}{\lambda^2}.
\end{equation}
Using \eqref{eq:Rsplit},
\begin{equation}
\int_b^L(L-t)R(t)dt=J(-1/2)-\sum_{k=1}^\infty J(a_k).
\end{equation}
The required infinite sums are bounded by
\begin{align}
\frac1{2a_N}
&\leq\sum_{k=N}^\infty a_k^{-2}
\leq\frac1{2a_N}+a_N^{-2},\label{eq:tailpower}\\
0&\leq\sum_{k=N}^\infty\frac{e^{-a_kL}}{a_k^2}
\leq\frac{e^{-a_NL}}{a_N^2(1-e^{-2L})},\label{eq:tailexp}\\
0&\leq\sum_{k=N}^\infty J(a_k)
\leq\frac{(L-b)e^{-a_Nb}}{a_N(1-e^{-2b})}.\label{eq:tailJ}
\end{align}
Equation~\eqref{eq:tailpower} is the integral bound for a decreasing
inverse-square function. Equation~\eqref{eq:tailexp} follows by bounding
all denominators below by $a_N^2$ and summing the geometric numerator.
For \eqref{eq:tailJ}, bound $L-t\leq L-b$, extend the exponential
integral to infinity, and then sum its geometric majorant.

For rational $x\geq0$, the checker encloses $e^x$ by its Taylor
polynomial through degree $m=160$. If $x<m+2$, the remaining positive
tail is at most
\begin{equation}
\frac{x^{m+1}}{(m+1)!}\frac1{1-x/(m+2)}.
\end{equation}
Negative arguments are handled by reciprocal intervals. All arguments
in this calculation obey the required bound.
For rational $x\geq1$, put $z=(x-1)/(x+1)$. Then
\begin{equation}
\log x=2\sum_{k=0}^{n-1}\frac{z^{2k+1}}{2k+1}+r_n,
\quad 0\leq r_n\leq
\frac{2z^{2n+1}}{(2n+1)(1-z^2)},
\end{equation}
with $n=48$. The alternating arctangent series, with 24 terms, and
Machin's formula
$\pi=16\arctan(1/5)-4\arctan(1/239)$ give rational endpoints for $\pi$.
The logarithm endpoints are then evaluated by monotonicity. Finally,
the elementary integral-test bounds
\begin{equation}
H_{128}-\log128-1/128<\gamma_E<H_{128}-\log128,
\qquad \log128=7\log2,
\end{equation}
are sufficient; no high-precision tabulated value of $\gamma_E$ is used.

Combining these intervals with \eqref{eq:constantgamma} and
\eqref{eq:U} produces an enclosure for $\mathcal U$ with approximate
endpoints
\begin{equation}
[-0.0494803686403,\,-0.0409467581007].
\end{equation}
The actual comparison in the checker is the exact rational assertion
\begin{equation}
\mathcal U_{\mathrm{upper}}<-1/25.
\end{equation}
The displayed decimals are for interpretation only. All additions,
multiplications, divisions, and interval endpoint comparisons in this
certificate use exact fractions. This is a reproducible computational
proof of a sign in the stated normalization, subject to review of the
short checker and the elementary bounds above.

\section{Rational certificates for the third investigation}
\label{app:round3}

The file \texttt{checks/round3/check\_round3.py} reuses the elementary
fraction interval operations and exponential, logarithm, and arctangent
enclosures of the second checker. Logarithm arguments are first reduced
by powers of two to $[1,2]$, so the logarithm series parameter has absolute
value at most $1/3$. The same $H_{128}$ bounds for $\gamma_E$ are used.
The program rejects optimized Python mode, since its imported helper
contains assertions. Its reported decimal values are explanatory; all
certifying comparisons are between fractions.

\subsection{The entire odd potential interval}

For \eqref{eq:oddscalar}, the checker evaluates interval extensions on
each of the 1980 closed cells
\[
[k/2000,(k+1)/2000],\qquad k=0,\ldots,1979.
\]
Their union is $[0,99/100]$. Every interval lower endpoint exceeds $1/100$;
the smallest is approximately $0.0176841$. On $[99/100,1)$ it bounds
$-\log(1-z)$ below by $\log100$, and evaluates the remaining rational
function on the closed interval $[99/100,1]$. This yields a lower bound
approximately $0.547915$. The denominator $1-4z/5$ is bounded below by
$1/5$ everywhere. The inequality $\cosh^2(7/40)<33/32$ is also verified
by fractions. These are interval covers of the full continuum, not
samples at grid points. The proof of the kernel estimate and the reduction
to this scalar function are in Theorem~\ref{thm:oddcompletion}.

\subsection{The linear-input signs}

For the unit linear input at length $L$, the contribution of mass $m>0$
to its kinetic jump form is
\begin{equation}
I_m(L)=\frac6{Lm^2}-\frac{24}{L^3m^4}
+e^{-mL}\left(\frac6{Lm^2}+\frac{24}{L^2m^3}
                         +\frac{24}{L^3m^4}\right).
\end{equation}
This follows by integrating
$2e^{-mt}(3t/L-2t^3/L^3)$ from zero to $L$ and adding the full-output
tail $2e^{-mL}/m$. Thus
\begin{equation}
\Qgam_{0,L}[f]=w_0+\sum_{k=0}^\infty I_{a_k}(L)
-\frac{24}{L^3}\bigl[2L\cosh(L/4)-8\sinh(L/4)\bigr]^2.
\end{equation}
The checker sets $L=1$, sums $k=0,\ldots,63$, and writes $m=a_{64}$.
The two inverse-power tails satisfy
\[
\frac1{2m}\leq\sum_{k=64}^\infty a_k^{-2}\leq\frac1{2m}+m^{-2},
\qquad
\frac1{6m^3}\leq\sum_{k=64}^\infty a_k^{-4}
\leq\frac1{6m^3}+m^{-4}.
\]
The positive exponential tail is at most
\[
\frac{e^{-mL}}{1-e^{-2L}}
\left(\frac6{Lm^2}+\frac{24}{L^2m^3}+\frac{24}{L^3m^4}\right).
\]
Combining these signed bounds proves the gamma enclosure. The arithmetic
coefficient is enclosed as $c=(\log2)e^{-(\log2)/2}$; the autocorrelation
and diagonal debit in Proposition~\ref{prop:primestabilizes} then give the
remaining enclosures with no square-root oracle or zero data.

\subsection{Separate identity diagnostics and controls}

At two quadrature orders and three lengths, the checker compares the
full-output autocorrelation formula against the folded cubic factor on a
complex odd polynomial. It integrates the endpoint logarithm by exact
polynomial moment formulas before numerical quadrature of the smooth
remainder. The largest observed absolute discrepancy is below $6\cdot10^{-16}$
for the unnormalized test inputs. These are floating-point identity checks,
not the sign certificate. The cap-reflection identity is also checked on
complex vectors. Positive symmetric and negative antisymmetric cap inputs
are both tested, preventing a mistaken assertion that the prime correction
is nonnegative on the entire odd sector.

\section{Verification, provenance, and remaining review}
\label{app:verification}

The supplied \texttt{checks/} directory contains the two original
diagnostic programs, the third investigation checker, and recorded outputs. They use Python and
NumPy and require no zeta-zero data. The second program's rational
certificate uses only standard-library fraction arithmetic. Its other
calculations, and all calculations in the first program, are separate
floating-point diagnostics.

The first program checks the digamma/jump normalization, the full-output
difference identity on complex polynomial inputs at two quadrature
orders, and the first-prime projection algebra. The second checks the
mass expansion at six frequencies against an independent digamma
evaluation, the auxiliary minimization, the coherent model and its
graded algebra, and the constant-input integrals at two quadrature
orders. The observed identity discrepancies are consistent with the
stated numerical truncations and floating-point precision. These
diagnostics do not certify an extension of the small-window theorem.

The normalization was taken from the project manuscripts
\cite{BakerFirst,BakerWeil} at repository commit
\texttt{c6eb19cc47b972b410b077be5af4e86a112afd95}.
The researcher's previously reported Fourier checks were accepted as
context; their large archive outputs were not required for the present
derivations. The data archive is named \texttt{szp-archive}.
Version 0.1 was prepared outside the source repository. Version 0.2 adds
the third investigation on branch \texttt{susy-positivity}, starting from
commit \texttt{4de517e}. No external archive or synced project document was
modified by this continuation.

The introductory history is based on the researcher's account and the
project brief. The mathematical results consolidate the two initial exploratory
investigations and the third continuation dated 11 September 2026. AI assistance was used in their derivation,
checking, and exposition, and in preparing this manuscript. No external
author was contacted and no independent review is implied.

Before this could be treated as a submission draft, the priority review
should compare the small-window argument with existing Dirichlet-form
methods and the pairwise restriction with established comparison-matrix
and factor-width results. The domains of the infinite tower and the
continuum pairwise definition also merit specialist review. The present
paper makes no general continuum Schur-limit theorem and no claim to
have constructed an interacting superspace theory.

One formula-convention question concerning negative time in a published
norm-map construction was recorded during the first investigation.
It is not needed for any proof here and is retained in the companion
research notes for further checking, rather than presented as an
established correction to that work.

\begin{thebibliography}{9}
\raggedright

\bibitem{BakerFirst}
E. Baker, \emph{Archimedean first-slab positivity for shifted zeta canonical systems}, project working
manuscript, 2026. Repository source:
\href{../../../shifted-zeta/first-slab-positivity/first_slab_positivity.tex}{First-slab positivity manuscript source}.
The source file, rather than a public publication, is the cited object.

\bibitem{BakerWeil}
E. Baker, \emph{Finite-horizon Weil coercivity and shifted-zeta contraction},
working draft, version 0.4, 10 September 2026. Repository source:
\href{../../../shifted-zeta/weil-depth/manuscript/finite_horizon_weil.tex}{Weil-depth manuscript source}.

\bibitem{DLMF}
NIST Digital Library of Mathematical Functions, Chapter 5,
especially \S\S5.7, 5.9, and 5.11.
\url{https://dlmf.nist.gov/5}.

\bibitem{KwasnickiMucha}
M. Kwa\'snicki and J. Mucha,
\emph{Extension technique for complete Bernstein functions of the Laplace
operator}, Journal of Evolution Equations (2018).
\href{https://doi.org/10.1007/s00028-018-0444-4}{doi:10.1007/s00028-018-0444-4}.
\href{https://arxiv.org/abs/1707.02475}{arXiv:1707.02475}.

\bibitem{Singleton}
A. Singleton, \emph{Superconformal quantum mechanics and the exterior
algebra}, 2014. \href{https://arxiv.org/abs/1403.4933}{arXiv:1403.4933}.

\bibitem{SuzukiHilbert}
M. Suzuki, \emph{On the Hilbert space derived from the Weil distribution},
2025. \href{https://doi.org/10.4153/S0008414X25101739}{doi:10.4153/S0008414X25101739}.
\href{https://arxiv.org/abs/2301.00421v3}{arXiv:2301.00421v3}.

\bibitem{SuzukiScrew}
M. Suzuki, \emph{Weil's quadratic form via the screw function}, 2026.
\href{https://arxiv.org/abs/2606.09096v1}{arXiv:2606.09096v1}.
This is the version consulted during the exploratory investigation;
a later revision exists and remains to be compared in the priority audit.

\bibitem{Yafaev}
D. R. Yafaev, \emph{On semibounded Wiener--Hopf operators}, 2016.
\href{https://arxiv.org/abs/1606.01361}{arXiv:1606.01361}.

\end{thebibliography}
\end{document}
